\documentclass[10pt,a4paper]{amsart}
\usepackage[margin=19mm]{geometry}
\usepackage{amsmath,amssymb,mathtools}
\usepackage[scr=rsfs,cal=euler]{mathalfa}
\usepackage{xfrac}
\usepackage[unicode,hidelinks,bookmarks=false]{hyperref}
\newcommand{\ie}{i.e.\ }
\newcommand{\cf}{cf.\ }
\newcommand{\resp}{respectively\ }
\newcommand{\Subsection}{\subsection}
\providecommand{\nobreakdash}{\nobreak\hbox{-}}
\setlength{\emergencystretch}{2em}
\multlinegap0pt
\usepackage{amssymb}
\usepackage{algorithm}\usepackage{algpseudocode}
\usepackage{float}
\newtheorem{proposition}{Proposition}
\newtheorem{theorem}{Theorem}
\title{Recursive algorithms for computing Birkhoff interpolation polynomials}
\author{Xue Jiang, Yuanhe Li, Zhe Li}
\date{Source: arXiv:2511.09014v2 (CC BY 4.0)}
\begin{document}
\maketitle

\noindent\emph{Source and licence.} Recursive algorithms for computing Birkhoff interpolation polynomials. Version arXiv:2511.09014v2 (CC BY 4.0), distributed by arXiv under the Creative Commons Attribution 4.0 International licence, http://creativecommons.org/licenses/by/4.0/, as recorded in the arXiv licence metadata for this exact version and shown on the abstract page https://arxiv.org/abs/2511.09014v2. Full licence notice: \texttt{licenses/arxiv-2511.09014-CC-BY-4.0.txt}.\par\smallskip

\noindent\textit{Editorial scope.} Counted unit: the article's theorem theorem (source0.tex lines 824--826). The article's complete original proof of it is delivered with it (source0.tex lines 827--899, one \texttt{proof} environment of 73 lines). No mathematics is written, repaired or re-derived: the statement and the proof are the article's own lines, in the article's own notation, and no retained line is substituted or re-flowed.  Retained uncounted as the source-local context the proof needs: lines 281--283; lines 508--516; lines 520--530; lines 592--599; lines 718--723.  Nothing was omitted from any retained span, so the package carries no omission record.  No retained line contains a citation, so the wrapper carries no bibliography and no bibliography entry.  No retained line contains a figure environment, an image-inclusion command or drawing code, so no image is delivered and the package declares no asset dependency.  Editorial formatting only, in addition: an AMS \texttt{amsart} preamble that restates the article's own declarations and macros used by the retained lines, namely the environments \texttt{proposition}, \texttt{theorem} on separate counters, together with the standard packages those lines need.  The retained environments print in this document as Proposition 1, Theorem 1 in order of appearance, whereas the article prints the same items with the numbering of its own full document.  The complete native source of the article as distributed in the arXiv e-print for this exact version is bundled unchanged as \texttt{sources/188944/source0.tex}.
\begin{align}\label{equa2.1}
   \delta_{x_{\beta_1}}\circ D_x^{\alpha_1},\delta_{x_{\beta_2}}\circ D_x^{\alpha_2},\dots,\delta_{x_{\beta_N}}\circ D_x^{\alpha_N},
\end{align}

\begin{equation}\label{equa3.3}
g_{j,k-1}(x):= \left(\left[\begin{array}{c|cccc}
            q_{k-1}(x) & q_0(x) & q_1(x) & \dots & q_{j}(x) \\ \hline
            q_{k-1}^{(\alpha_1)}(x_{\beta_1}) & q_0^{(\alpha_1)}(x_{\beta_1}) & q_1^{(\alpha_1)}(x_{\beta_1}) & \dots & q_{j}^{(\alpha_1)}(x_{\beta_1}) \\
            q_{k-1}^{(\alpha_2)}(x_{\beta_2}) & q_0^{(\alpha_2)}(x_{\beta_2}) & q_1^{(\alpha_2)}(x_{\beta_2}) & \dots & q_{j}^{(\alpha_2)}(x_{\beta_2}) \\
            \vdots & \vdots & \vdots & ~ & \vdots \\
            q_{k-1}^{(\alpha_{j+1})}(x_{\beta_{j+1}}) & q_0^{(\alpha_{j+1})}(x_{\beta_{j+1}}) & q_1^{(\alpha_{j+1})}(x_{\beta_{j+1}}) & \dots & q_{j}^{(\alpha_{j+1})}(x_{\beta_{j+1}}) \\
            \end{array}\right]/V_{j}\right).
\end{equation}

\begin{proposition}\label{prop:V.V}
Let $\textup{span}\{q_0(x), q_1(x),\dots,q_{N-1}(x)\}$ be a strongly proper interpolation space for the Birkhoff interpolation problem $(\ref{equa2.1})$. For any fixed $k=1,2,\dots,N$, let $g_{-1,k-1}(x)=q_{k-1}(x)$. Then  
\begin{align*}
    g_{k-2,k-1}^{(\alpha_i)}(x_{\beta_i})=\left\{
    \begin{array}{ll}
      0, & i=1,2,\dots,k-1; \\
      \dfrac{|V_{k-2}|}{|V_{k-1}|}\neq 0, & i=k.
    \end{array}
  \right.
\end{align*}
\end{proposition}

\begin{theorem}\label{TH7}
Let $\textup{span}\{q_0(x), q_1(x),\dots,q_{N-1}(x)\}$ be a strongly proper interpolation space for the Birkhoff interpolation problem $(\ref{equa2.1})$.
$g_{j,k-1}(x)$, $k=1,2,\dots,N, j=0,1,\dots,k-2$ are defined by $(\ref{equa3.3})$.
Let $p_{-1}(x)=0$, then the interpolation polynomial %$k=2, 3, \dots, N$, we have
\begin{equation*}
p_{k-1}(x)=p_{k-2}(x)+\frac{y_{\beta_k,\alpha_k}-p_{k-2}^{(\alpha_k)}(x_{\beta_k})}{g_{k-2,k-1}^{(\alpha_k)}(x_{\beta_k})}g_{k-2,k-1}(x),~~k=1, 2, \dots, N.
\end{equation*}
\end{theorem}

\begin{theorem}\label{gx} 
 Under the same assumptions as Theorem $\ref{TH7}$, for any fixed $k=1,2,\dots,N$,
\begin{equation*}
g_{j,k-1}(x)=g_{j-1,k-1}(x)-\frac{g_{j-1,k-1}^{(\alpha_{j+1})}(x_{\beta_{j+1}})}{g_{j-1,j}^{(\alpha_{j+1})}(x_{\beta_{j+1}})}g_{j-1,j}(x),~ j=0, 1, \dots, k-2.
\end{equation*}
\end{theorem}

\begin{theorem}
 The output  $\{g_{-1,0}(x), g_{0,1}(x), \dots, g_{N-2,N-1}(x)\}$ is a Newton-type   basis for the Birkhoff interpolation problem $(\ref{equa2.1})$. Furthermore, it is  a set of strongly  proper interpolation basis, and $p_{N-1}(x)$ is the corresponding interpolation polynomial.
\end{theorem}

\begin{proof}
From Proposition \ref{prop:V.V} and Theorem \ref{gx}, to prove that
\begin{displaymath}
  \{g_{-1,0}(x), g_{0,1}(x), \dots, g_{N-2,N-1}(x)\}  
\end{displaymath}
is a Newton-type basis,
we only need to prove that the recursively updated
\begin{displaymath}
  \{g_{-1,0}(x), g_{-1,1}(x), \dots, g_{-1,N-1}(x)\}  
\end{displaymath}
 is a strongly proper basis.

We will use induction on the number of interpolation conditions.
When $k=1$, $ \delta_{x_{\beta_1}}\circ D_x^{\alpha_1}(g_{-1,0}(x))= \delta_{x_{\beta_1}}\circ D_x^{\alpha_1}(x^{\alpha_1})=\alpha_1!\neq 0.$
When $k=2$, $g_{-1,1}(x)=x^{\alpha_1+1+s}, s=1,2,\dots$. This corresponds to the execution process of $k=2$ in Algorithm 1. 

First, assume $g_{-1,1}(x)=x^{\alpha_1+1}$ holds and compute $g_{0,1}(x)$ using line 16 of Algorithm 1. If $g_{0,1}(x)$ satisfies the judgment condition (line 18 in Algorithm 1), this computational step is completed and
$g_{-1,1}(x),g_{0,1}(x)$ are stored. Otherwise, the degree of $g_{-1,1}(x)$ is increased, and 
$g_{0,1}(x)$ is recomputed recursively until a 
$g_{0,1}(x)$
 that satisfies 
$g_{0,1}^{(\alpha_2)}(x_{\beta_2})\neq 0$
 is obtained.
It is straightforward to verify that 
\begin{displaymath}
    g_{0,1}^{(\alpha_1)}(x_{\beta_1})=g_{-1,1}^{(\alpha_1)}(x_{\beta_1})-\frac{g_{-1,1}^{(\alpha_1)}(x_{\beta_1})}{g_{-1,0}^{(\alpha_1)}(x_{\beta_1})}g_{-1,0}^{(\alpha_1)}(x_{\beta_1})=0.
\end{displaymath} 
Thus, $\{g_{-1,0}(x),g_{0,1}(x)\}$ forms a set of strongly proper Newton-type basis for the first two interpolation conditions. Since $g_{0,1}(x)$ is a linear combination of $g_{-1,1}(x)$ and $g_{-1,0}(x)$, 
$\{g_{-1,0}(x),g_{-1,1}(x)\}$ constitutes a set of strongly proper monomial basis for the first two conditions.


Now assume the result has been proved for the first  $k-1$ interpolation conditions, that is,
$\{g_{-1,0}(x), g_{-1,1}(x), \dots, g_{-1,k-2}(x)\}$ is a strongly proper basis. By Proposition \ref{prop:V.V}, we know that $\{g_{-1,0}(x), g_{0,1}(x), \dots, g_{k-3,k-2}(x)\}$ is a Newton-type basis. 
First, we prove that $\{g_{-1,0}(x), g_{0,1}(x), \dots, g_{k-2,k-1}(x)\}$ is a Newton-type basis for the first $k$ conditions. By the induction hypothesis, it suffices to further verify that 
\begin{displaymath}
  g_{k-2,k-1}(x)=g_{k-3,k-1}(x)-\frac{g_{k-3,k-1}^{(\alpha_{k-1})}(x_{\beta_{k-1}})}{g_{k-3,k-2}^{(\alpha_{k-1})}(x_{\beta_{k-1}})}g_{k-3,k-2}(x) , 
\end{displaymath}
 satisfies: 
$\textup{(\romannumeral 1) }g_{k-2,k-1}^{(\alpha_k)}(x_{\beta_k})\neq 0;$ 
$\textup{(\romannumeral 2) }g_{k-2,k-1}^{(\alpha_i)}(x_{\beta_i})=0, \forall i=1,2,\dots,k-1$.

By line 18 of Algorithm 1, we know that the first inequality holds.
Note that $g_{-1,k-1}(x), g_{0,k-1}(x), \dots, g_{k-2,k-1}(x)$ are polynomials of the same degree. 
For $i=0,1,\dots,k-2$, it is easy to verify that $g_{i,k-1}(x)$ vanishes at the first $i+1$ conditions, i.e.,
\begin{displaymath}
    g_{i,k-1}^{(\alpha_j)}(x_{\beta_j})=0, \forall j=1,2,\dots,i+1.
\end{displaymath}
Furthermore, 
%\begin{displaymath}
%    g_{k-2,k-1}(x)=g_{k-3,k-1}(x)-\frac{g_{k-3,k-1}^{(\alpha_{k-1})}(x_{\beta_{k-1}})}{g_{k-3,k-2}^{(\alpha_{k-1})}(x_{\beta_{k-1}})}g_{k-3,k-2}(x)  , 
%\end{displaymath} 
%and 
\begin{displaymath}
    g_{k-2,k-1}^{(\alpha_i)}(x_{\beta_i})=g_{k-3,k-1}^{(\alpha_i)}(x_{\beta_i})-\frac{g_{k-3,k-1}^{(\alpha_{k-1})}(x_{\beta_{k-1}})}{g_{k-3,k-2}^{(\alpha_{k-1})}(x_{\beta_{k-1}})}g_{k-3,k-2}^{(\alpha_i)}(x_{\beta_i}),  \forall i=1,2,\dots,k-1.
\end{displaymath}
It is obvious that $g_{k-2,k-1}^{(\alpha_i)}(x_{\beta_i})=0$ when $i=k-1$. For $i=1,2,\dots,k-2$, as shown in the preceding proof, we know that 
$g_{k-3,k-1}^{(\alpha_i)}(x_{\beta_i})=0$, and by the induction hypothesis, \[
\left\{
\begin{aligned}
&g_{k-3,k-2}^{(\alpha_i)}(x_{\beta_i})=0, \\
&g_{k-3,k-2}^{(\alpha_{k-1})}(x_{\beta_{k-1}})\neq 0.
\end{aligned}
\right.
\] 
Hence
\begin{displaymath}
    g_{k-2,k-1}^{(\alpha_i)}(x_{\beta_i})=0, \forall i=1,2,\dots,k-1.
\end{displaymath} Thus we have proved that $\{g_{-1,0}(x), g_{0,1}(x), \dots, g_{k-2,k-1}(x)\}$ forms a set of strongly proper Newton-type basis.

Since $$\textup{span}\{g_{-1,0}(x), g_{-1,1}(x), \dots, g_{-1,k-1}(x)\}=\textup{span}\{g_{-1,0}(x),g_{0,1}(x),\dots,g_{k-2,k-1}(x)\},$$ 
it follows that
$\{g_{-1,0}(x), g_{-1,1}(x), \dots, g_{-1,k-1}(x)\}$ constitutes a set of strongly proper monomial basis.
\end{proof}

\end{document}
